\documentclass[11pt]{article}
\usepackage{aa}

\begin{document}
\zestaw{3}{Funkcje tworzące}

\begin{teoria}[Jak czytać ten dokument]
Każde zadanie ma: przepisaną \textbf{treść} (szare pudełko), potrzebną
\textbf{teorię} (niebieskie), \textbf{intuicję} (żółte), \textbf{rozwiązanie
formalne} i~\textbf{odpowiedź} (zielone). Wyniki sprawdzono skryptem
\texttt{weryfikacja/aa03.py} (współczynniki szeregów w~ułamkach, pełne
przeglądy ciągów, dokładne momenty rozkładów).
Teoria idzie za wykładem 3 („Szeregi potęgowe i~funkcje tworzące”);
rzeczy spoza wykładu są oznaczone jako uzupełnienie.
Notacja probabilistyczna jak na wykładzie 2: $\Pr\{A\}$, $E[X]$.
\end{teoria}

% =====================================================================
\section{Teoria z wykładu: funkcje tworzące}

\begin{definicja}[Funkcja tworząca]
Funkcją tworzącą ciągu $\langle a_n\rangle=a_0,a_1,\dots$ nazywamy szereg
potęgowy
\[
G(z)=\sum_{n=0}^\infty a_nz^n.
\]
W~zadaniach zmienna nazywa się zwykle $x$, a~funkcja $A(x)$, $F(x)$ itp.
Piszemy $[x^n]A(x)=a_n$ (współczynnik przy $x^n$).
Dla ciągu dwuindeksowego $a_{n,m}$ analogicznie (uzupełnienie):
\[
A(x,y)=\sum_{n,m}a_{n,m}x^ny^m.
\]
\end{definicja}

\begin{intuicja}
Funkcja tworząca to „wieszak” na cały ciąg: ciąg staje się \emph{jednym
obiektem}, na którym można robić algebrę (dodawać, mnożyć, różniczkować),
a~potem odczytać współczynniki. Żeby poprawnie traktować $G(z)$ jako
\emph{funkcję} dla $z\in S$, szereg powinien być zbieżny dla $z\in S$.
\textbf{W~zastosowaniach} (jak na wykładzie) nie dbamy o~zbieżność ---
do rozwiązania dochodzimy formalnie, manipulując na szeregach, \textbf{ale}
na końcu potwierdzamy wynik, np.\ indukcyjnie.
\end{intuicja}

\begin{fakt}[Przekształcenia funkcji tworzących, z~wykładu]\label{fakt:slownik}
Dla $A(x)=\sum a_nx^n$, $B(x)=\sum b_nx^n$ (wszystko formalnie):
\begin{align*}
pA+qB&\;\leftrightarrow\;pa_n+qb_n,\\
x^mA(x)&\;\leftrightarrow\;a_{n-m}\quad(\text{przesunięcie; } a_{<0}=0),\\
A(x)B(x)&\;\leftrightarrow\;\sum_{k=0}^na_kb_{n-k}\quad(\text{iloczyn Cauchy'ego}).
\end{align*}
Mnożenie przez liczbę to szczególny przypadek iloczynu Cauchy'ego. Inny
szczególny przypadek: $b_n=1$ dla wszystkich $n$, $B(x)=\frac1{1-x}$, daje
\textbf{sumy częściowe}:
\[
\frac{A(x)}{1-x}=a_0+(a_0+a_1)x+(a_0+a_1+a_2)x^2+\cdots
\]
\emph{Uzupełnienie} --- różniczkowanie i~całkowanie:
\begin{align*}
A'(x)&\;\leftrightarrow\;(n+1)\,a_{n+1},\\
xA'(x)&\;\leftrightarrow\;n\,a_n,\\
\int_0^xA(t)\,dt&\;\leftrightarrow\;\frac{a_{n-1}}{n}\quad(n\ge1).
\end{align*}
\end{fakt}

\begin{fakt}[Podstawowe szeregi]\label{fakt:szeregi}
\begin{align*}
\frac1{1-x}&=\sum_{n\ge0}x^n,&
\frac1{1-ax}&=\sum_{n\ge0}a^nx^n,\\
\frac1{(1-x)^2}&=\sum_{n\ge0}(n+1)x^n,&
(1+x)^m&=\sum_{k}\binom mkx^k.
\end{align*}
Pierwszy wzór jak na wykładzie: $G(x)=\sum_{n\ge0}x^n$ i~$xG(x)=\sum_{n\ge1}x^n$,
więc $(1-x)G(x)=1$. Drugi to pierwszy po podstawieniu $x\to ax$, trzeci to pochodna
pierwszego: $\bigl(\frac1{1-x}\bigr)'=\frac1{(1-x)^2}=\sum nx^{n-1}$, czwarty
to wzór dwumianowy.
\end{fakt}

\subsection*{Przykład z wykładu: wzór zwarty na \texorpdfstring{$F_n$}{F(n)}}

$F_0=0$, $F_1=1$, $F_n=F_{n-1}+F_{n-2}$ dla $n\ge2$. Niech
$G(z)=\sum_{n\ge0}F_nz^n$. Wtedy $zG(z)$ i~$z^2G(z)$ to funkcje tworzące
ciągów $\langle F_{n-1}\rangle$ i~$\langle F_{n-2}\rangle$, więc
$(1-z-z^2)G(z)$ to funkcja tworząca ciągu $\langle F_n-F_{n-1}-F_{n-2}\rangle$.
Wyrazy od $z^2$ wzwyż znikają z~rekurencji:
\begin{align*}
(1-z-z^2)G(z)&=F_0+(F_1-F_0)z+0\cdot z^2+\dots=z,\\
G(z)&=\frac z{1-z-z^2}.
\end{align*}
Mianownik ma pierwiastki $\frac{-1\pm\sqrt5}2$, czyli
$1-z-z^2=(1-\varphi z)(1-\hat\varphi z)$, gdzie
\[
\varphi=\frac{1+\sqrt5}2,\qquad\hat\varphi=1-\varphi=\frac{1-\sqrt5}2.
\]
Rozkład na ułamki proste i~$\frac1{1-\varphi z}=1+\varphi z+\varphi^2z^2+\cdots$ dają
\begin{align*}
G(z)&=\frac1{\sqrt5}\Bigl(\frac1{1-\varphi z}-\frac1{1-\hat\varphi z}\Bigr)
=\sum_{n\ge0}\frac{\varphi^n-\hat\varphi^n}{\sqrt5}z^n,\\
F_n&=\frac{\varphi^n-\hat\varphi^n}{\sqrt5}\qquad\text{(wzór Bineta)}.
\end{align*}
Ten sam schemat (rekurencja $\to$ równanie na $G$ $\to$ ułamki proste $\to$
odczytanie współczynników) stosujemy w~zad.~4 i~9.

\subsection*{Funkcja tworząca prawdopodobieństwa}

\begin{definicja}[z wykładu]\label{def:pgf}
Niech $X_n$ będzie zmienną losową o~wartościach w~$\{0,1,2,\dots\}$, np.\
liczbą kroków algorytmu na losowych danych rozmiaru $n$, i~niech
$p_{nk}=\Pr\{X_n=k\}$. \emph{Funkcja tworząca (rozkładu)
prawdopodobieństwa} zmiennej $X_n$ to
\[
G_n(z)=\sum_{k=0}^\infty p_{nk}z^k.
\]
\end{definicja}

\begin{fakt}[z wykładu]\label{fakt:pgf}
$G_n(1)=1$ (suma prawdopodobieństw) oraz
\begin{align*}
E[X_n]&=G_n'(1),\\
\operatorname{Var}[X_n]&=G_n''(1)+G_n'(1)-G_n'(1)^2.
\end{align*}
\end{fakt}
\begin{proof}
$G_n'(z)=\sum_{k\ge1}kp_{nk}z^{k-1}$, więc $G_n'(1)=\sum_kkp_{nk}=E[X_n]$.
Dalej $G_n''(1)=\sum_kk(k-1)p_{nk}$, a~z~definicji
$\operatorname{Var}[X_n]=\sum_kp_{nk}(k-E[X_n])^2$:
\begin{align*}
\operatorname{Var}[X_n]&=\sum_kk^2p_{nk}-2G_n'(1)\sum_kkp_{nk}+G_n'(1)^2\sum_kp_{nk}\\
&=\sum_kk(k-1)p_{nk}+\sum_kkp_{nk}-G_n'(1)^2\\
&=G_n''(1)+G_n'(1)-G_n'(1)^2.\qedhere
\end{align*}
\end{proof}

Na wykładzie wprowadzono też \emph{formalne operatory} dla dowolnego
$G(z)=\sum_ka_kz^k$:
\[
E(G)=\sum_{k\ge1}k\,a_k,\qquad
\operatorname{Var}(G)=\sum_{k\ge0}\bigl(k-E(G)\bigr)^2a_k.
\]
Jeśli $G$ jest funkcją tworzącą prawdopodobieństwa zmiennej $X$, to
oczywiście $E[X]=E(G)$ i~$\operatorname{Var}[X]=\operatorname{Var}(G)$.

\begin{lemat}[z wykładu]\label{lem:iloczyn}
Jeśli $G$, $H$ są funkcjami tworzącymi ciągów $\langle a_n\rangle$,
$\langle b_n\rangle$ oraz $\sum a_i=\sum b_i=1$, to dla iloczynu
Cauchy'ego $G\cdot H$:
\begin{align*}
E(G\cdot H)&=E(G)+E(H),\\
\operatorname{Var}(G\cdot H)&=\operatorname{Var}(G)+\operatorname{Var}(H).
\end{align*}
\end{lemat}
\begin{proof}[Dowód (na wykładzie pominięty)]
Z~Faktu~\ref{fakt:pgf}, zastosowanego formalnie, $E(G)=G'(1)$ i~$\operatorname{Var}(G)=G''(1)+G'(1)-G'(1)^2$
(dowód nie korzystał z~tego, że $a_k$ to prawdopodobieństwa, tylko z~$G(1)=1$).
Przy $G(1)=H(1)=1$:
\begin{align*}
(GH)'(1)&=G'(1)+H'(1),\\
(GH)''(1)&=G''(1)+2G'(1)H'(1)+H''(1),
\end{align*}
więc
\begin{align*}
\operatorname{Var}(GH)&=G''+2G'H'+H''+G'+H'-(G'+H')^2\\
&=\bigl(G''+G'-G'^2\bigr)+\bigl(H''+H'-H'^2\bigr)
\end{align*}
(wszystko w~punkcie 1).
\end{proof}

\begin{intuicja}
Iloczyn funkcji tworzących prawdopodobieństwa niezależnych zmiennych $X$, $Y$
to funkcja tworząca $X+Y$ (bo $\Pr\{X+Y=n\}=\sum_k\Pr\{X=k\}\Pr\{Y=n-k\}$ ---
dokładnie iloczyn Cauchy'ego). Lemat mówi więc znane: przy niezależności
wartości oczekiwane i~wariancje się dodają.
\end{intuicja}

\subsection*{Przykład z wykładu: szukanie maksimum}

Problem sekretarki: znajdź maksimum w~$a[1..n]$.
\begin{algorithm}[H]
\caption{\textsc{Max} (z wykładu)}
\begin{algorithmic}[1]
\State $\mathit{maxi}\gets1$
\For{$i\gets2$ \textbf{do} $n$}
  \If{$a[i]>a[\mathit{maxi}]$}
     \State $\mathit{maxi}\gets i$ \Comment{$(*)$}
  \EndIf
\EndFor
\State \Return $a[\mathit{maxi}]$
\end{algorithmic}
\end{algorithm}
(To \textsc{Zatrudnij} z~zestawu~2, tyle że start od $\mathit{maxi}=1$,
więc $X_n$ to liczba zatrudnień minus 1.) Dane to losowa permutacja $a_1,\dots,a_n$ zbioru
$[n]$; $X_n$ --- liczba wykonań instrukcji $(*)$.

\emph{Rekurencja na $p_{nk}$.} Niech $c_{nk}$ to liczba permutacji z~$X_n=k$,
czyli $p_{nk}=c_{nk}/n!$. Jeśli $a_n=n$, to $X_n$ jest o~1 większe niż dla
$a_1,\dots,a_{n-1}$; jeśli $a_n\ne n$ ($n-1$ wyborów $a_n$), to takie samo.
Zatem
\begin{align*}
c_{nk}&=c_{n-1,k-1}+(n-1)\,c_{n-1,k},\\
p_{nk}&=\frac1n\,p_{n-1,k-1}+\frac{n-1}n\,p_{n-1,k},
\end{align*}
oraz $p_{1k}=I\{k=0\}$ (dla $n=1$ pętla się nie wykonuje).

\emph{Funkcja tworząca.} $G_1(z)=1$, a~dla $n>1$, mnożąc rekurencję przez
$z^k$ i~sumując (pierwszy składnik to przesunięcie o~1):
\[
G_n(z)=\frac zn\,G_{n-1}(z)+\frac{n-1}n\,G_{n-1}(z)=\frac{z+n-1}n\,G_{n-1}(z).
\]
($G_n$ to wielomian, bo $p_{nk}=0$ dla $k\ge n$.)

\emph{Wartość oczekiwana.} Pochodna iloczynu i~$G_{n-1}(1)=1$:
\[
G_n'(1)=\frac1n\cdot1+1\cdot G_{n-1}'(1),\qquad G_1'(1)=0,
\]
więc $E[X_n]=G_n'(1)=\sum_{k=2}^n\frac1k=H_n-1$ ($-1$, bo iterujemy od $i=2$).

\emph{Wariancja.} Rozwijając rekurencję,
\[
G_n(z)=\frac{z+n-1}n\cdot\frac{z+n-2}{n-1}\cdots\frac{z+1}2\cdot G_1(z)
=\prod_{k=2}^nQ_k(z),
\]
gdzie $Q_k(z)=\frac{z+k-1}k$. Każde $Q_k$ ma sumę współczynników $Q_k(1)=1$,
$Q_k'=\frac1k$, $Q_k''=0$, więc $E(Q_k)=\frac1k$,
$\operatorname{Var}(Q_k)=0+\frac1k-\frac1{k^2}$. Z~Lematu~\ref{lem:iloczyn}:
\begin{align*}
E[X_n]&=\sum_{k=2}^n\frac1k=H_n-1,\\
\operatorname{Var}[X_n]&=\sum_{k=2}^n\Bigl(\frac1k-\frac1{k^2}\Bigr)
=H_n-H_n^{(2)},
\end{align*}
gdzie $H_n^{(2)}=\sum_{k=1}^n\frac1{k^2}\to\frac{\pi^2}6$. Zatem
$\operatorname{Var}[X_n]=\ln n+\gamma-\frac{\pi^2}6+o(1)$ ($\gamma\approx0{,}577$
--- stała Eulera).
\emph{Sens:} $Q_k$ to funkcja tworząca indykatora „$a_k$ jest maksimum
spośród $a_1,\dots,a_k$” (prawd.\ $\frac1k$); te indykatory są niezależne
i~$X_n$ jest ich sumą.

% =====================================================================
\section{Zadanie 1. Funkcja tworząca ciągu $1/n$}
\begin{tresc}
Wyznacz zwartą postać funkcji tworzącej
$F(x)=\sum_{n=1}^{\infty}\frac1nx^n$.
\end{tresc}

\begin{intuicja}
Współczynnik $\frac1n$ to „wynik całkowania $x^{n-1}$”:
$\int_0^xt^{n-1}dt=\frac{x^n}{n}$. Czyli $F$ to całka z~najprostszego szeregu
$1+x+x^2+\cdots=\frac1{1-x}$. Całka z~$\frac1{1-t}$ to logarytm.
\end{intuicja}

\formalnie
Różniczkując wyraz po wyrazie (Fakt~\ref{fakt:slownik}),
\[
F'(x)=\sum_{n\ge1}x^{n-1}=\frac1{1-x},\qquad F(0)=0.
\]
Stąd, dla $|x|<1$,
\[
F(x)=\int_0^x\frac{dt}{1-t}=-\ln(1-x)=\ln\frac1{1-x}.
\]
(Kontrola: $-\ln(1-x)=x+\frac{x^2}2+\frac{x^3}3+\cdots$ --- szereg Maclaurina
logarytmu.)
\begin{odpowiedz}
$F(x)=\displaystyle\sum_{n\ge1}\frac{x^n}{n}=\ln\frac1{1-x}$ \quad dla $|x|<1$.
\end{odpowiedz}

% =====================================================================
\section{Zadanie 2. Liczby harmoniczne}
\begin{tresc}
Korzystając z~poprzedniego zadania, wykaż, że funkcja
$H(x)=\frac1{1-x}\ln\frac1{1-x}$ jest funkcją tworzącą ciągu liczb
harmonicznych: $H_n=\sum_{k=1}^n\frac1k$. Następnie wyznacz funkcję
tworzącą dla ciągu $\bigl(H_{n-1}/n\bigr)_{n\ge1}$.
\end{tresc}

\begin{intuicja}
$H_n=\frac11+\frac12+\dots+\frac1n$ to \emph{suma częściowa} ciągu $\frac1k$,
a~mnożenie przez $\frac1{1-x}$ właśnie sumuje częściowo
(Fakt~\ref{fakt:slownik}). Ciąg $\frac{H_{n-1}}n$ to z~kolei „$H$ po
całkowaniu” (dzielenie przez $n$ i~przesunięcie).
\end{intuicja}

\formalnie
\emph{Część 1.} Z~zad.~1: $\ln\frac1{1-x}=\sum_{k\ge1}\frac1kx^k$ (czyli
$b_0=0$, $b_k=\frac1k$). Mnożymy przez $\frac1{1-x}=\sum_{j\ge0}x^j$; ze splotu
\[
[x^n]H(x)=\sum_{k=0}^nb_k\cdot1=\sum_{k=1}^n\frac1k=H_n.
\]

\emph{Część 2.} Niech $a_n=H_n$ dla $n\ge0$ ($H_0=0$). Całkując $H(x)$ od $0$
do $x$ (Fakt~\ref{fakt:slownik}):
\[
\int_0^xH(t)\,dt=\sum_{n\ge1}\frac{H_{n-1}}{n}x^n.
\]
Lewą stronę liczymy podstawieniem $u=-\ln(1-t)$, $du=\frac{dt}{1-t}$:
\[
\int_0^x\frac{-\ln(1-t)}{1-t}\,dt=\int_0^{-\ln(1-x)}u\,du=\frac12\ln^2(1-x).
\]
(Szybka kontrola przez pochodną: $\bigl(\frac12\ln^2(1-x)\bigr)'=
\ln(1-x)\cdot\frac{-1}{1-x}=H(x)$ \checkmark.)
\begin{odpowiedz}
$H(x)=\dfrac1{1-x}\ln\dfrac1{1-x}=\sum_{n\ge0}H_nx^n$ oraz
\[
\sum_{n\ge1}\frac{H_{n-1}}{n}x^n=\frac12\ln^2\frac1{1-x}.
\]
\end{odpowiedz}
\emph{Ciekawostka.} $\frac{H_{n-1}}n$ to dokładnie prawdopodobieństwo
„dokładnie dwóch zatrudnień” z~zadania~2 zestawu~2. Teraz mamy jego funkcję
tworzącą w~zwartej postaci.

% =====================================================================
\section{Zadanie 3. Funkcja tworząca dwóch zmiennych}
\begin{tresc}
Wyznacz zwartą postać funkcji tworzącej
\[
F(x,y)=\sum_{n=0}^{\infty}\sum_{m=0}^{n}m\binom nm x^ny^m.
\]
\end{tresc}

\begin{intuicja}
Gdyby nie czynnik $m$, wewnętrzna suma byłaby dwumianem Newtona
$\sum_m\binom nmy^m=(1+y)^n$. Czynnik $m$ przy $y^m$ robi operator
$y\frac{\partial}{\partial y}$ (jak $xA'(x)$ w~słowniku). Potem zostaje zwykły
szereg geometryczny po $n$.
\end{intuicja}

\formalnie
Dla ustalonego $n$ (Fakt~\ref{fakt:szeregi}): $\sum_{m=0}^n\binom nmy^m=(1+y)^n$.
Zastosujmy do obu stron operator $y\frac{\partial}{\partial y}$:
\[
\sum_{m=0}^nm\binom nmy^m=y\cdot n(1+y)^{n-1}.
\]
Zatem, sumując po $n$ i~podstawiając $u=x(1+y)$,
\begin{align*}
F(x,y)&=\sum_{n\ge1}n\,y(1+y)^{n-1}x^n
=xy\sum_{n\ge1}n\bigl(x(1+y)\bigr)^{n-1}\\
&=xy\sum_{n\ge1}nu^{n-1}=\frac{xy}{(1-u)^2}
=\frac{xy}{(1-x-xy)^2}.
\end{align*}
(Szereg $\sum nu^{n-1}=\frac1{(1-u)^2}$ --- Fakt~\ref{fakt:szeregi}; zbieżność dla $|x|(1+|y|)<1$, ale jako
szereg formalny wzór jest prawdziwy zawsze.)

\emph{Kontrola:} dla $n=2$ współczynnik przy $x^2$ to
$\sum_m m\binom2my^m=2y+2y^2$; ze wzoru
$xy\cdot\bigl(1+2x(1+y)+\dots\bigr)$ dostajemy przy $x^2$: $xy\cdot2x(1+y)
=2x^2y+2x^2y^2$ \checkmark.
\begin{odpowiedz}
$F(x,y)=\dfrac{xy}{(1-x-xy)^2}$.
\end{odpowiedz}

% =====================================================================
\section{Zadanie 4. Liczby Lucasa}
\begin{tresc}
Wykorzystując funkcje tworzące, wykaż, że
\[
\sum_{k=0}^{\lfloor n/2\rfloor}\frac{n}{n-k}\binom{n-k}{k}=L_n,
\]
gdzie $L_n$ to $n$-ta liczba Lucasa ($L_0=2$, $L_1=1$,
$L_{n+2}=L_{n+1}+L_n$).
\end{tresc}

\begin{teoria}[Funkcja tworząca liczb Lucasa]
Tak jak dla $F_n$ w~przykładzie z~wykładu:
niech $L(x)=\sum_{n\ge0}L_nx^n$. Z~rekurencji $L_n=L_{n-1}+L_{n-2}$ dla
$n\ge2$:
\[
L(x)-xL(x)-x^2L(x)=L_0+(L_1-L_0)x=2-x,
\]
bo wszystkie wyrazy od $x^2$ wzwyż się skracają. Zatem
\[
L(x)=\frac{2-x}{1-x-x^2},\qquad
L(x)-2=\frac{x+2x^2}{1-x-x^2}=\sum_{n\ge1}L_nx^n.
\]
\end{teoria}

\begin{intuicja}
Zamiast udowadniać tożsamość indukcją (która wymaga zgadnięcia, co
wzmocnić), liczymy funkcję tworzącą lewej strony i~widzimy, że to ta sama
funkcja wymierna. Trzy sztuczki: (1)~sumę po $k$ zamieniamy na sumę po
parze $(m,k)$, gdzie $m=n-k$, czyli $x^n=x^{m+k}$; (2)~$\sum_k\binom mkx^k=(1+x)^m$;
(3)~czynnik $\frac nm$ rozbijamy na „$1+\frac km$” i~korzystamy z~tożsamości
$\frac km\binom mk=\binom{m-1}{k-1}$.
\end{intuicja}

\formalnie
Niech $b_n$ oznacza lewą stronę dla $n\ge1$ (dla $n=0$ wyraz $\frac00\binom00$
jest nieokreślony; tożsamość czytamy dla $n\ge1$). Podstawmy $m=n-k$.
Wtedy $n=m+k$ i~dla $m\ge1$
\begin{align*}
\frac{n}{n-k}\binom{n-k}{k}&=\frac{m+k}{m}\binom mk=\binom mk+\frac km\binom mk\\
&=\binom mk+\binom{m-1}{k-1}.
\end{align*}
(Sumujemy po $k\le\lfloor n/2\rfloor$, czyli $k\le m$ --- dokładnie te $k$,
dla których $\binom mk\ne0$; analogicznie dla $\binom{m-1}{k-1}$ i~$k\ge1$.)
Każda para $(m,k)$ z~$m\ge1$, $0\le k\le m$ pojawia się dokładnie raz
(przy $n=m+k$; dla ustalonego $n$ to skończenie wiele par), więc
\begin{align*}
\sum_{n\ge1}b_nx^n
&=\sum_{m\ge1}\sum_{k}\binom mkx^{m+k}\\
&\quad+\sum_{m\ge1}\sum_{k\ge1}\binom{m-1}{k-1}x^{m+k}.
\end{align*}
Pierwsza suma:
\begin{align*}
\sum_{m\ge1}x^m\sum_k\binom mkx^k&=\sum_{m\ge1}x^m(1+x)^m\\
&=\frac1{1-x(1+x)}-1=\frac1{1-x-x^2}-1.
\end{align*}
W~drugiej podstawiamy $m'=m-1$, $k'=k-1$ ($x^{m+k}=x^2x^{m'+k'}$):
\begin{align*}
x^2\sum_{m'\ge0}x^{m'}\sum_{k'}\binom{m'}{k'}x^{k'}&=x^2\sum_{m'\ge0}x^{m'}(1+x)^{m'}\\
&=\frac{x^2}{1-x-x^2}.
\end{align*}
Razem
\begin{align*}
\sum_{n\ge1}b_nx^n&=\frac{1+x^2}{1-x-x^2}-1\\
&=\frac{x+2x^2}{1-x-x^2}=\sum_{n\ge1}L_nx^n,
\end{align*}
więc $b_n=L_n$ dla $n\ge1$. \qed

\emph{Kontrola:} $n=4$: $k=0$: $\frac44\binom40=1$; $k=1$: $\frac43\binom31=4$;
$k=2$: $\frac42\binom22=2$; suma $7=L_4$ \checkmark.
\begin{odpowiedz}
$\displaystyle\sum_{k}\frac n{n-k}\binom{n-k}k$ ma funkcję tworzącą
$\dfrac{x+2x^2}{1-x-x^2}$, a~to jest funkcja tworząca $(L_n)_{n\ge1}$.
(Pomocniczo: $\sum_k\binom{n-k}{k}=F_{n+1}$ ma funkcję tworzącą
$\frac1{1-x-x^2}$ --- liczby Fibonacciego.)
\end{odpowiedz}

% =====================================================================
\section{Zadanie 5. Funkcja tworząca momentów}
\begin{tresc}
Funkcją tworzącą momentów zmiennej losowej $X\colon\Omega\to\{x_0,x_1,\ldots\}
\subseteq\R$ nazywamy szereg $G(t)=\sum_{n=0}^\infty\Pr\{X=x_n\}e^{tx_n}$.
Korzystając z~funkcji $G(t)$, wyznacz wartość oczekiwaną oraz wariancję $X$.
\end{tresc}

\begin{teoria}[Wariancja]
$\operatorname{Var}X=E\bigl[(X-E[X])^2\bigr]$ --- średni kwadrat odchylenia
od średniej. Z~liniowości wartości oczekiwanej (zestaw~2),
dla $\mu=E[X]$:
\begin{align*}
\operatorname{Var}X&=E[X^2-2\mu X+\mu^2]\\
&=E[X^2]-2\mu E[X]+\mu^2=E[X^2]-\mu^2.
\end{align*}
\end{teoria}

\begin{intuicja}
$G(t)=E[e^{tX}]$ (średnia z~funkcji $e^{tX}$ --- wprost z~definicji $E$ jako
$\sum x\Pr\{X=x\}$). Rozwińmy $e^{tX}=1+tX+\frac{t^2X^2}{2}+\cdots$ i~weźmy
średnią:
\[
G(t)=1+E[X]\,t+E[X^2]\frac{t^2}{2!}+E[X^3]\frac{t^3}{3!}+\cdots
\]
Moment $E[X^k]$ siedzi więc przy $\frac{t^k}{k!}$ --- stąd nazwa „funkcja
tworząca momentów”. Żeby go wyciągnąć, $k$ razy różniczkujemy i~wstawiamy
$t=0$ (wszystkie wyższe wyrazy znikają).
\end{intuicja}

\formalnie
Zakładamy, że $G(t)$ jest skończona w~pewnym otoczeniu zera (dla zmiennych
o~skończonym zbiorze wartości --- zawsze). Wtedy szereg
$\sum_n\Pr\{X=x_n\}e^{tx_n}$ można różniczkować wyraz po wyrazie
(szereg pochodnych jest zbieżny jednostajnie w~mniejszym otoczeniu zera,
bo $|x_n|^ke^{t x_n}$ szacuje się przez $e^{\delta|x_n|}\le
e^{\delta x_n}+e^{-\delta x_n}$ z~dokładnością do stałej). Mamy
\begin{align*}
G'(t)&=\sum_n\Pr\{X=x_n\}\,x_ne^{tx_n},\\
G''(t)&=\sum_n\Pr\{X=x_n\}\,x_n^2e^{tx_n}.
\end{align*}
Wstawiając $t=0$ ($e^0=1$):
\begin{align*}
G'(0)&=\sum_n x_n\Pr\{X=x_n\}=E[X],\\
G''(0)&=\sum_n x_n^2\Pr\{X=x_n\}=E[X^2].
\end{align*}
Z~wzoru na wariancję
\[
\operatorname{Var}X=G''(0)-\bigl(G'(0)\bigr)^2.
\]
Przy okazji: $G(0)=\sum_n\Pr\{X=x_n\}=1$.

Wariant: $(\ln G)'(t)=\frac{G'}{G}$ i~$(\ln G)''=\frac{G''G-G'^2}{G^2}$,
a~w~zerze (bo $G(0)=1$) dostajemy $(\ln G)'(0)=E[X]$ oraz
$(\ln G)''(0)=\operatorname{Var}X$. Często to wygodniejsze, bo logarytm
zamienia iloczyny na sumy.
\begin{odpowiedz}
$E[X]=G'(0)$, \quad $\operatorname{Var}X=G''(0)-G'(0)^2=(\ln G)''(0)$.
\end{odpowiedz}

% =====================================================================
\section{Zadanie 6. Rozkład jednostajny}
\begin{tresc}
Zmienna losowa $X$ ma rozkład jednostajny: $\Pr\{X=j\}=\frac1n$ dla
$j=1,2,\ldots,n$. Przedstaw postać zwartą funkcji tworzącej
prawdopodobieństwa $F(z)$ oraz funkcji tworzącej momentów $G(t)$ zmiennej
losowej $X$. Ponadto wyznacz wartość oczekiwaną oraz wariancję zmiennej $X$.
\end{tresc}

\begin{teoria}[Funkcja tworząca prawdopodobieństwa]\label{teoria:pgf}
To funkcja z~Definicji~\ref{def:pgf} (na wykładzie $G_n(z)$, w~treści
zadań $F(z)$, bo $G(t)$ to tu funkcja tworząca momentów):
$F(z)=\sum_{j\ge0}\Pr\{X=j\}z^j=E[z^X]$. Z~Faktu~\ref{fakt:pgf}
\[
E[X]=F'(1),\qquad
\operatorname{Var}X=F''(1)+F'(1)-F'(1)^2,
\]
bo $F'(1)=E[X]$ i~$F''(1)=E[X(X-1)]$ („silniowe” momenty).
Ponadto $G(t)=E[e^{tX}]=F(e^t)$, więc z~reguły łańcuchowej
$G'(0)=F'(1)$ i~$G''(0)=F''(1)+F'(1)$ --- to
te same liczby co w~zad.~5. (Pochodne w~$z=1$ biorą się jako pochodne
szeregu; dla nas wszystkie szeregi mają promień zbieżności $>1$ albo są
wielomianami.)
\end{teoria}

\begin{intuicja}
$F(z)$ to „wieszak na rozkład”: współczynnik przy $z^j$ to
prawdopodobieństwo wartości $j$. Dla rozkładu jednostajnego to po prostu
wielomian $\frac1n(z+z^2+\dots+z^n)$, czyli suma geometryczna.
\end{intuicja}

\formalnie
\begin{align*}
F(z)&=\frac1n\sum_{j=1}^nz^j=\frac{z\,(z^n-1)}{n\,(z-1)}\qquad(z\ne1),\\
G(t)&=F(e^t)=\frac1n\sum_{j=1}^ne^{jt}=\frac{e^t\bigl(e^{nt}-1\bigr)}{n\,(e^t-1)}\qquad(t\ne0),
\end{align*}
oczywiście $F(1)=G(0)=1$. Pochodne w~$z=1$ (lub $t=0$) z~postaci
ułamkowej dają wyrażenie $\frac00$, więc rachunek wygodniej prowadzić na
postaci sumy:
\begin{align*}
F'(1)&=\frac1n\sum_{j=1}^nj=\frac{n+1}{2},\\
F''(1)&=\frac1n\sum_{j=1}^nj(j-1)=\frac1n\cdot\frac{(n-1)n(n+1)}{3}=\frac{n^2-1}{3},
\end{align*}
bo $\sum j(j-1)=\sum j^2-\sum j=\frac{n(n+1)(2n+1)}6-\frac{n(n+1)}2
=\frac{n(n+1)(n-1)}3$. Stąd
\begin{align*}
\operatorname{Var}X&=\frac{n^2-1}{3}+\frac{n+1}2-\frac{(n+1)^2}4\\
&=\frac{4n^2-4+6n+6-3n^2-6n-3}{12}=\frac{n^2-1}{12}.
\end{align*}
\begin{odpowiedz}
\begin{align*}
F(z)&=\frac{z(z^n-1)}{n(z-1)},&
G(t)&=\frac{e^t(e^{nt}-1)}{n(e^t-1)},\\
E[X]&=\frac{n+1}2,&
\operatorname{Var}X&=\frac{n^2-1}{12}.
\end{align*}
\end{odpowiedz}

% =====================================================================
\section{Zadanie 7. Rozkład dwumianowy}
\begin{tresc}
Zmienna losowa $X$ ma rozkład dwumianowy: dla ustalonego $p\in(0,1)$,
$\Pr\{X=j\}=\binom njp^j(1-p)^{n-j}$ dla $j=0,1,\ldots,n$. Przedstaw postać
zwartą funkcji tworzącej prawdopodobieństwa $F(z)$ oraz funkcji tworzącej
momentów $G(t)$ zmiennej losowej $X$. Ponadto wyznacz wartość oczekiwaną
oraz wariancję zmiennej $X$.
\end{tresc}

\begin{intuicja}
Prawdopodobieństwa to wyrazy wzoru dwumianowego dla $(q+pz)^n$, gdzie
$q=1-p$ --- funkcja tworząca sama się „zwija”.
Zmienna $X$ to liczba sukcesów w~$n$ niezależnych próbach, a~$z^X=\prod_iz^{X_i}$ dla
indykatorów prób, więc $F(z)=(q+pz)^n$ to $n$-ta potęga funkcji pojedynczej
próby (tu: niezależność zamienia sumę zmiennych na iloczyn funkcji
tworzących --- intuicja do Lematu~\ref{lem:iloczyn}).
\end{intuicja}

\formalnie
Niech $q=1-p$. Ze wzoru dwumianowego
\[(a+b)^n=\sum_j\binom nja^jb^{n-j}\]
dla $a=pz$, $b=q$:
\begin{align*}
F(z)&=\sum_{j=0}^n\binom nj(pz)^jq^{n-j}=(q+pz)^n,\\
G(t)&=F(e^t)=\bigl(q+pe^t\bigr)^n.
\end{align*}
Pochodne $F$:
\begin{align*}
F'(z)&=np\,(q+pz)^{n-1},\\
F''(z)&=n(n-1)p^2(q+pz)^{n-2}.
\end{align*}
W~$z=1$ mamy $q+p=1$, więc $F'(1)=np$, $F''(1)=n(n-1)p^2$ i
\begin{align*}
\operatorname{Var}X&=n(n-1)p^2+np-n^2p^2\\
&=np-np^2=np(1-p).
\end{align*}
\emph{Drugi sposób} (przez $\ln G$, zad.~5): $\ln G=n\ln(q+pe^t)$,
\begin{align*}
(\ln G)'&=\frac{npe^t}{q+pe^t},\\
(\ln G)''&=n\,\frac{pe^t(q+pe^t)-p^2e^{2t}}{(q+pe^t)^2}\\
&=\frac{npqe^t}{(q+pe^t)^2}.
\end{align*}
Przy $t=0$ ($q+p=1$) dostajemy $E[X]=np$, $\operatorname{Var}X=npq$ \checkmark.

\emph{Trzeci sposób} (Lemat~\ref{lem:iloczyn} z~wykładu): $F=Q^n$ dla
$Q(z)=q+pz$, a~$Q(1)=1$, $Q'=p$, $Q''=0$, więc $E(Q)=p$,
$\operatorname{Var}(Q)=0+p-p^2=pq$. Iloczyn $n$ kopii daje
$E[X]=np$ i~$\operatorname{Var}X=npq$ \checkmark.
\begin{odpowiedz}
\begin{align*}
F(z)&=(1-p+pz)^n,&
G(t)&=(1-p+pe^t)^n,\\
E[X]&=np,&
\operatorname{Var}X&=np(1-p).
\end{align*}
\end{odpowiedz}

% =====================================================================
\section{Zadanie 8. Rozkład geometryczny}
\begin{tresc}
Zmienna losowa $X$ ma rozkład geometryczny: dla ustalonego $p\in(0,1)$,
$\Pr\{Y=j\}=(1-p)^{j-1}p$ dla $j=1,\ldots,n$. Przedstaw postać zwartą
funkcji tworzącej prawdopodobieństwa $F(z)$ oraz funkcji tworzącej
momentów $G(t)$ zmiennej losowej $X$. Ponadto wyznacz wartość oczekiwaną
oraz wariancję zmiennej $X$.
\end{tresc}

\begin{uwaga}[o~treści]
W~treści jest literówka: raz $X$, raz $Y$ --- chodzi o~tę samą zmienną.
Zakres $j=1,\dots,n$ też nie pasuje do rozkładu geometrycznego: suma
$\sum_{j=1}^n(1-p)^{j-1}p=1-(1-p)^n<1$, więc to nie byłby rozkład
prawdopodobieństwa. Przyjmuję standardowe $j=1,2,3,\dots$ (liczba prób do
pierwszego sukcesu, jak w~zestawie~2, fakt o~rozkładzie geometrycznym); wtedy
$\sum_{j\ge1}(1-p)^{j-1}p=1$.
\end{uwaga}

\begin{intuicja}
Prawdopodobieństwa tworzą ciąg geometryczny o~ilorazie $q=1-p$, a~funkcja
tworząca ciągu geometrycznego to $\frac{1}{1-qz}$. Tylko trzeba doliczyć
czynnik $pz$ (pierwsza próba kończy się sukcesem, dlatego $j$ zaczyna się
od 1).
\end{intuicja}

\formalnie
Niech $q=1-p$. Ze wzoru na sumę szeregu geometrycznego, dla $|qz|<1$:
\begin{align*}
F(z)&=\sum_{j\ge1}q^{j-1}p\,z^j=pz\sum_{i\ge0}(qz)^i=\frac{pz}{1-qz},\\
G(t)&=F(e^t)=\frac{pe^t}{1-qe^t}\qquad\bigl(t<\ln\tfrac1q\bigr).
\end{align*}
Promień zbieżności $F$ to $\frac1q>1$, więc pochodne w~$z=1$ istnieją:
\begin{align*}
F'(z)&=\frac{p(1-qz)+pz\cdot q}{(1-qz)^2}=\frac{p}{(1-qz)^2},\\
F''(z)&=\frac{2pq}{(1-qz)^3}.
\end{align*}
Ponieważ $1-q=p$:
\[
F'(1)=\frac p{p^2}=\frac1p,\qquad F''(1)=\frac{2pq}{p^3}=\frac{2q}{p^2},
\]
\begin{align*}
\operatorname{Var}X&=\frac{2q}{p^2}+\frac1p-\frac1{p^2}
=\frac{2q+p-1}{p^2}\\
&=\frac{2-2p+p-1}{p^2}=\frac{1-p}{p^2}.
\end{align*}
Zgadza się to z~zestawem~2, gdzie $E[X]=\frac1p$ wyszło z~sumy ogonów.

\emph{Drugi sposób} (przez $\ln G$): $\ln G=\ln p+t-\ln(1-qe^t)$,
\[
(\ln G)'=1+\frac{qe^t}{1-qe^t},\qquad
(\ln G)''=\frac{qe^t}{(1-qe^t)^2}.
\]
Dla $t=0$: $E[X]=1+\frac qp=\frac1p$ i~$\operatorname{Var}X=\frac q{p^2}$ \checkmark.
\begin{odpowiedz}
\begin{align*}
F(z)&=\frac{pz}{1-(1-p)z},&
G(t)&=\frac{pe^t}{1-(1-p)e^t},\\
E[X]&=\frac1p,&
\operatorname{Var}X&=\frac{1-p}{p^2}.
\end{align*}
\end{odpowiedz}

% =====================================================================
\section{Zadanie 9. Ciągi jedynek i dwójek}
\begin{tresc}
Niech $G_n$ będzie liczbą ciągów o~wyrazach należących do zbioru $\{1,2\}$,
których suma jest równa $n$ (oczywiście $G_n=F_{n+1}$, gdzie $F_n$ jest
$n$-tą liczbą Fibonacciego). Dla $m\ge0$ niech $G_{n,m}$ oznacza liczbę
takich ciągów zawierających dokładnie $m$ jedynek. Pokaż, jak wykorzystać
funkcję tworzącą dwóch zmiennych
\[
G(x,y)=\sum_{n,m\ge0}G_{n,m}x^ny^m
\]
do wyznaczenia oczekiwanej liczby jedynek w~ciągu o~wyrazach sumujących się
do $n$. Wyznacz tę liczbę jako funkcję zmiennej $n$ z~dokładnością do
$\O(1)$.
\end{tresc}

\begin{intuicja}
\textbf{Krok 1: funkcja tworząca.} Ciąg to dowolnie długie słowo z~liter
„1” i~„2”. Litera „1” wnosi do sumy 1 i~jedną jedynkę, więc ma wagę $xy$;
litera „2” wnosi 2 do sumy i~żadnej jedynki: waga $x^2$. Ciąg to
ciąg liter, więc jego waga to \emph{iloczyn} wag liter, a~wszystkie ciągi
razem to $\sum_{k\ge0}(xy+x^2)^k=\frac1{1-xy-x^2}$ (suma po długości $k$).

\textbf{Krok 2: średnia z~jednej funkcji.} „Losowy ciąg o~sumie $n$” to
losowy wybór jednego z~$G_n$ ciągów, każdy z~prawd.\ $\frac1{G_n}$. Liczba
jedynek $M$ ma wtedy
\[
\Pr\{M=m\}=\frac{G_{n,m}}{G_n},\qquad
E[M]=\frac{\sum_mmG_{n,m}}{G_n}.
\]
Licznik to współczynnik przy $x^n$
w~$\frac{\partial G}{\partial y}$ przy $y=1$ (pochodna po $y$ mnoży przez
$m$ --- jak w~zad.~3 --- a~$y=1$ sumuje po $m$), mianownik --- w~$G(x,1)$.
To dokładnie wzór $E[X]=F'(1)$ z~zad.~6, tylko dla
funkcji tworzącej prawdopodobieństwa $\frac{[x^n]G(x,y)}{G_n}$ zmiennej $M$.
\end{intuicja}

\formalnie
\emph{Funkcja tworząca.} Ciągowi o~$m$ jedynkach i~$j$ dwójkach (suma
$n=m+2j$) przypisujemy jednomian $x^ny^m=(xy)^m(x^2)^j$. Po rozwinięciu
$(xy+x^2)^k$ każdy składnik odpowiada wyborowi jednej litery na każdej
z~$k$ pozycji, czyli dokładnie jednemu ciągowi długości $k$, więc
\[
G(x,y)=\sum_{k\ge0}(xy+x^2)^k=\frac1{1-xy-x^2}.
\]
Podstawienie $y=1$: $G(x,1)=\frac1{1-x-x^2}=\sum_nG_nx^n$, czyli
$G_n=F_{n+1}$ (jak w~treści i~zad.~4: $F_0=0$, $F_1=F_2=1$).

\emph{Wartość oczekiwana.} Niech $h_n=\sum_mmG_{n,m}$ (łączna liczba jedynek
we wszystkich ciągach o~sumie $n$), tak że $E[M]=h_n/G_n$. Pochodna po $y$
mnoży współczynnik przy $y^m$ przez $m$, a~$y=1$ sumuje po $m$:
\begin{align*}
\frac{\partial G}{\partial y}(x,y)&=\frac{x}{(1-xy-x^2)^2},\\
D(x):=\frac{\partial G}{\partial y}(x,1)&=\frac{x}{(1-x-x^2)^2}=\sum_nh_nx^n.
\end{align*}
Zostaje odczytać współczynniki $D$. Oznaczmy $Q=1-x-x^2$ i~$\Phi=\frac1Q=\sum_n
F_{n+1}x^n$. Kwadrat w~mianowniku przypomina pochodną: z~Faktu~\ref{fakt:slownik}
\[
\Phi'(x)=\frac{1+2x}{Q^2}=\sum_{n\ge0}(n+1)F_{n+2}\,x^n.
\]
Trzeba więc zapisać $\frac x{Q^2}$ przez $\Phi'$ i~$\Phi$. Pomaga
\emph{dopełnienie do kwadratu}: $Q=\frac54-\bigl(x+\frac12\bigr)^2$, czyli
\[
(1+2x)^2+4Q=5.
\]
Stąd
\[
\frac1{Q^2}=\frac{(1+2x)^2+4Q}{5Q^2}=\frac{1+2x}5\,\Phi'+\frac45\Phi,
\]
a~ponieważ $x=\frac{(1+2x)-1}2$,
\begin{align*}
D=\frac x{Q^2}&=\frac12\,\Phi'-\frac12\cdot\frac1{Q^2}\\
&=\frac12\Phi'-\frac{1+2x}{10}\Phi'-\frac25\Phi\\
&=\frac{2-x}5\,\Phi'-\frac25\,\Phi.
\end{align*}
Odczytujemy współczynnik przy $x^n$ ($x\Phi'$ to przesunięcie: $[x^n]\,x\Phi'=nF_{n+1}$):
\begin{align*}
h_n&=\frac{2(n+1)F_{n+2}-nF_{n+1}-2F_{n+1}}5\\
&=\frac{2(n+1)F_{n+2}-(n+2)F_{n+1}}5.
\end{align*}
(Kontrola: $h_1=\frac{4\cdot2-3}5=1$, $h_2=\frac{6\cdot3-4\cdot2}5=2$,
$h_3=\frac{8\cdot5-5\cdot3}5=5$ --- np.\ dla $n=3$ ciągi $(1,1,1),(1,2),(2,1)$
mają razem $3+1+1=5$ jedynek \checkmark.) Ponieważ $L_{n+1}=F_n+F_{n+2}=
2F_{n+2}-F_{n+1}$, to samo można zapisać jako
$h_n=\frac{(n+1)L_{n+1}-F_{n+1}}5$. Zatem
\[
E[M]=\frac{h_n}{F_{n+1}}=\frac15\Bigl(2(n+1)\frac{F_{n+2}}{F_{n+1}}-(n+2)\Bigr).
\]

\emph{Asymptotyka.} Niech $\varphi=\frac{1+\sqrt5}2$, $\psi=\hat\varphi=\frac{1-\sqrt5}2$
(pierwiastki $t^2=t+1$, $|\psi|<1$). Wzór Bineta $F_k=\frac{\varphi^k-\psi^k}{\sqrt5}$
(przykład z~wykładu) daje
\[
F_{k+1}-\varphi F_k=\frac{\psi^k(\varphi-\psi)}{\sqrt5}=\psi^k,
\]
więc $\frac{F_{n+2}}{F_{n+1}}=\varphi+\frac{\psi^{n+1}}{F_{n+1}}
=\varphi+\O(\varphi^{-2n})$ (bo $F_{n+1}=\Theta(\varphi^n)$, $|\psi|=\varphi^{-1}$).
Wstawiając i~korzystając z~$2\varphi-1=\sqrt5$:
\begin{align*}
E[M]&=\frac{(2\varphi-1)(n+1)-1}5+\O\bigl(n\varphi^{-2n}\bigr)\\
&=\frac{n+1}{\sqrt5}-\frac15+o(1)
=\frac n{\sqrt5}+\O(1).
\end{align*}
Stała to $\frac1{\sqrt5}-\frac15\approx0{,}247$. Przykłady: $n=2$: ciągi
$(1,1),(2)$, średnio $1$; $n=3$: średnio $\frac53$ --- oba zgodne ze wzorem
dokładnym.

\emph{Sens:} $\frac1{\sqrt5}\approx0{,}447$, czyli w~typowym ciągu jest
ok.\ $0{,}447n$ jedynek i~$0{,}276n$ dwójek (bo $m+2j=n$).
\begin{odpowiedz}
$G(x,y)=\dfrac1{1-xy-x^2}$ i
\begin{align*}
E[\#\text{jedynek}]&=\frac{[x^n]\,\partial_yG(x,1)}{[x^n]\,G(x,1)}\\
&=\frac{2(n+1)F_{n+2}-(n+2)F_{n+1}}{5F_{n+1}}\\
&=\frac{n}{\sqrt5}+\O(1).
\end{align*}
\end{odpowiedz}

\end{document}
